
\documentclass[sigconf]{acmart}
% ===== Don't touch this :) ======
\settopmatter{printacmref=false} % Removes citation information below abstract
\renewcommand\footnotetextcopyrightpermission[1]{} % removes footnote with conference information in first column
\renewcommand\footnotetextauthorsaddresses[1]{}
\pagestyle{plain} % removes running headers
% ================================

\usepackage{graphicx}

% for red coloring
\def\bad#1{\textcolor{red}{#1}}

\begin{document}
\title{CS561 Final Report: Implementing SIEVE in the Postgres Bufferpool}

\author{Ryan Crosier}
\affiliation{%
    \institution{Boston University}
    \country{}
}
\author{Will Garlington}
\affiliation{%
    \institution{Boston University}
    \country{}
}
\author{Jackson Gilstrap}
\affiliation{%
    \institution{Boston University}
    \country{}
}

\begin{abstract}
    The buffer cache is a key data structure in the DBMS that helps reduce unecessary disk I/Os. Eviction policies help determine when data in a cache is no longer needed and removes it, creating room for new data. Some of the most common algorithms are Least Recently Used (LRU), Least Frequently Used (LFU), First In First Out (FIFO), and many others. PostgreSQL uses the CLOCK \bad{citethis} algorithm, a modification of LRU designed to minimize the overhead of maintaining an ordered list of recently used pages. For our project, we implement SIEVE, an eviction policy originally designed as an LRU alternative for web caches, into PostgreSQL. SIEVE is advantageous due to its lock-less design and its ability to be integrated into existing cache architectures. Whilst SIEVE was originally made for the web, there could be possible use for it in database systems such as PostgreSQL. Our experiments benchmark Sieve against other eviction policies using \textit{pgbench}. We find that SIEVE provides little (if any) improvement in terms of throughput and cache-hit ratio over the default buffer cache eviction policy on mixed and skewed workloads.
\end{abstract}

\maketitle

% =============================================================================
\section{Introduction}
% =============================================================================

\subsection{Motivation}

The most successful cache eviction policy introduced in recent years is SIEVE \cite{SIEVE}, an algorithm which is even simpler to implement than LRU. While originally designed for web-caching workloads, it has been found to be useful in many workloads with a high degree of data skew. While the paper introducing SIEVE concedes that it may not perform well for buffer cache workloads due to the presence of range scans, many real-world workloads may have a much higher proportion of point queries for which the benefits of SIEVE can be exploited.

\subsection{Problem Statement}

The problem we are exploring is whether or not there are workloads for which SIEVE outperforms other simple eviction policies for the bufferpool in PostgreSQL. PostgreSQL uses the CLOCK algorithm for buffer cache eviction by default and we will also implement LRU as another point of comparison. We will then explore workloads with different data skews and determine whether SIEVE grants any performance benefits or reduces write amplification.

\subsection{Contributions}

We make the following contributions:

\begin{itemize}
    \item We implement a framework in PostgreSQL for implementing and switching between different buffer cache eviction policies at compile time.
    \item We present results indicating that SIEVE does not offer any significant performance benefits over CLOCK in skewed and mixed workloads.
\end{itemize}


% =============================================================================
\section{Background}
% =============================================================================
For decades cache eviction has been primarily dominated by very simple primitive algorithms such as First In First Out(FIFO) and Least Recently Used (LRU). 
More complex algorithms such as Adaptive Replacement Cache (ARC) typically tend to have better results, however require the developer to implement hundreds of lines of code. 
This persistent trade off between simplicity and efficiency has been continuously discussed and researched in system research. 
\cite{SIEVE}SIEVE is a recent advancement in this space, originally designed for and evaluated against web cache workloads. 
The authors categorize SIEVE as a cache eviction primitive similar to LRU or FIFO. Structurally, it is built upon a single FIFO queue and utilizes a "moving hand" pointer to track whether data has been recently accessed. 
This hand scans from the tail toward the head of the queue, resetting the "visited" bit of objects it encounters. An object is only evicted if its visited bit is 0; otherwise, it is retained in its original position, effectively "sifting" the cache to preserve popular data. 

% =============================================================================
\section{Architecture}
% =============================================================================

\subsection{Buffer Policy Framework}

To facilitate the testing of several eviction policies in a common 'environment', we've implemented an eviction pollicy framework on which all policies can be built. This framework is made up of two primary components: a BufferStrategyCommon which holds a buffer lock and other policy data, and an EvictionVtable which defines virtual functions to be implemented by new policies. Because \emph{freelist.c}, where eviction logic is implemented, does not dictate when or how buffers enter or leave the buffer pool, our virtual functions can focus purely on the mechanics of evictions and insertions.

\subsection{Sieve Implementation}

Our SIEVE implementation was designed with static arrays to minimize the overhead of memory reallocation. Two static arrays sized to \emph{NBuffers}, the number of buffers allowed in the pool, combine to recreate the behavior of a linked list. Each insertion or \emph{GetBuffer} (find victim) call is treated as an atomic action where a Spinlock is held until the function completes. After every insertion or deletion, the ID of the buffer is used to access the array and either mark the buffer as the new head (insertion) or designate the buffer as unreachable (eviction).

% ================================================
\section{Benchmarking}
% =============================================================================

\subsection{Hardware}

We perform all of our evaluations on a system running Linux Mint (a Debian fork). The system has an AMD Ryzen 5 5600X CPU with 6 Cores at 3.7GHz base clock, and 32 GB of DDR4 memory. The system is equipped with two storage devices: a 256 GB SATA SSD and a 1 TB NVMe SSD.

\subsection{Postgres Setup}

Postgres data is stored on the NVMe SSD for every test, and the WAL is stored on the SATA SSD. Benchmarks use the standard pgbench tables with auto-vacuuming disabled. The tables are vacuumed manually after each test to ensure a clean database for the next test.

\subsection{Pgbench Workloads}

We test over various workloads with different data localities and read/write ratios. For a skewed workload, we draw points from a Zipfian distribution with $\alpha=1.01$.

% =============================================================================
\section{Conclusion}
% =============================================================================

We have implemented and evaluated SIEVE against Clock Sweep for various synthetic workloads and found that there is little difference in throughput and cache hit ratio. Our next steps are to additionally measure the write amplification of different eviction policies over different workloads and determine if SIEVE is able to pull ahead in any of those scenarios. We would also like to test synthetic workloads include range scans, which should perform worse according to~\cite{SIEVE}.


% reference
{
    \bibliographystyle{ACM-Reference-Format}
    \bibliography{biblio}
} 

\end{document}

